\section*{Declaraciones}
\addcontentsline{toc}{section}{Declaraciones}

\subsection*{Disponibilidad de código}
El código fuente completo está disponible públicamente bajo licencia MIT en
\url{https://github.com/gleiston-guerrero/PFC-Presustentaciones-2026}, y archivado de forma permanente
en Zenodo con DOI \href{https://doi.org/10.5281/zenodo.22445216}{10.5281/zenodo.22445216} (versión
\texttt{v1.0.1}, el cierre real de esta Entrega Final). El DOI de concepto
\href{https://doi.org/10.5281/zenodo.21988563}{10.5281/zenodo.21988563} resuelve siempre a la
última versión archivada; la versión anterior \texttt{v1.0.0} permanece disponible por separado en
\href{https://doi.org/10.5281/zenodo.21988564}{10.5281/zenodo.21988564} (ver \texttt{docs/ZENODO.md}).

\subsection*{Disponibilidad de datos}
Los datos crudos de todas las mediciones empíricas citadas en este informe (corridas k6, reportes
Lighthouse, reportes OWASP ZAP, reporte find-sec-bugs, reporte JaCoCo) están versionados en el mismo
repositorio, con su procedencia documentada en \texttt{docs/mediciones/DATA-PROVENANCE.md}, y además
depositados por separado en Zenodo como un dataset independiente con licencia Creative Commons
Attribution 4.0 (CC-BY 4.0) y DOI propio
\href{https://doi.org/10.5281/zenodo.22398713}{10.5281/zenodo.22398713} (DOI de concepto
\texttt{10.5281/zenodo.22398712}; ver \texttt{docs/ZENODO-DATASET.md}), siguiendo el principio de
citación independiente entre software y datos. El instrumento SUS (Sección~\ref{sec:evaluacion}) sí
generó datos de participantes humanos: 15 respuestas anónimas (el formulario no pide nombre, solo el rol
dentro del sistema), de las cuales 4 tienen fecha de aplicación verificable y se usan en el resultado
reportado; las 11 restantes quedan versionadas como evidencia pero fuera del resultado de cierre hasta
confirmar su fecha (\texttt{docs/mediciones/sus/SUS-RESULTS.md}). El formulario llevaba una nota impresa
de consentimiento informado (participación voluntaria y anónima, con fines exclusivamente académicos),
pero no se recolectó una firma ni un consentimiento individual explícito por participante --- brecha
declarada, no subsanada retroactivamente.

\subsection*{Disponibilidad de materiales}
Los scripts de análisis (\texttt{scripts/gen-figuras.py}, \texttt{scripts/perf-analysis.ipynb},
\texttt{scripts/sus-analysis.ipynb}, \texttt{scripts/validate-traceability.sh},
\texttt{scripts/audit-sql-dynamic.sh}), el script de carga k6, y los archivos de configuración de
despliegue están en el mismo repositorio bajo la misma licencia.

\subsection*{Declaración ética}
Este proyecto no involucró experimentación con seres humanos con datos reales a la fecha de este informe
(el instrumento SUS existe pero no se ha aplicado). El protocolo de consentimiento informado para cuando
se aplique está documentado en \texttt{docs/etica/consentimientos/plantilla-consentimiento.md}. Se
declara explícitamente, como parte de esta misma sección, que una versión anterior del documento ético
del proyecto (\texttt{docs/etica/ETHICS.md}) afirmaba falsamente que existían 10 formularios de
consentimiento firmados en custodia física para una encuesta SUS que en realidad nunca se aplicó; esa
afirmación fue detectada y corregida durante este proceso, y se documenta aquí por transparencia.

\subsection*{Taxonomía CRediT (Contributor Roles Taxonomy)}

Se aplican los 14 roles estándar de CRediT \cite{brand2015} a los cuatro integrantes del equipo. La
Tabla~\ref{tab:credit} reemplaza la versión anterior de \texttt{CONTRIBUTORS.md}, que mezclaba roles
CRediT reales con descriptores de proyecto no estandarizados (p. ej. ``UI/UX Design'', ``DevOps
(Docker)''); esta tabla usa exclusivamente los 14 términos oficiales de la taxonomía.

\begin{table}[H]
\centering
\scriptsize
\begin{tabular}{lcccc}
\toprule
\textbf{Rol CRediT} & \textbf{Alava} & \textbf{Moncayo} & \textbf{Zamora} & \textbf{Barreto} \\
\midrule
Conceptualization & \checkmark & & & \\
Data curation & & & \checkmark & \\
Formal analysis & \checkmark & & & \checkmark \\
Funding acquisition & --- & --- & --- & --- \\
Investigation & \checkmark & \checkmark & \checkmark & \checkmark \\
Methodology & \checkmark & & \checkmark & \\
Project administration & \checkmark & & & \\
Resources & & & & \checkmark \\
Software & \checkmark & \checkmark & \checkmark & \checkmark \\
Supervision & --- & --- & --- & --- \\
Validation & & & \checkmark & \checkmark \\
Visualization & & \checkmark & & \\
Writing -- original draft & & & \checkmark & \\
Writing -- review \& editing & \checkmark & \checkmark & \checkmark & \checkmark \\
\bottomrule
\end{tabular}
\caption{Assignment of the 14 CRediT roles. ``---'' marks a role that genuinely does not apply to any
team member in this project (there was no external funding nor a formal supervisor beyond the
faculty director, whose role is declared separately, not as a co-author).}
\label{tab:credit}
\end{table}

\textit{Nota:} Alava Alvarado concentró la seguridad backend (JWT/Spring Security) y la administración
del proyecto; Moncayo Loor, el frontend Angular; Zamora Arias, el diseño de base de datos y la
documentación técnica; Barreto Rosado, las pruebas y la infraestructura Docker. Los cuatro participaron
en revisión de código y redacción/edición de la documentación (fila ``Writing -- review \& editing''),
consistente con lo ya declarado en \texttt{CONTRIBUTORS.md}.

\subsection*{Conflictos de interés}
Los autores declaran no tener conflictos de interés financieros ni personales que pudieran haber
influido en este trabajo.

\subsection*{Financiamiento}
Este proyecto no recibió financiamiento externo. Se desarrolló como parte de las actividades académicas
de la asignatura Aplicaciones Web, Quinto Nivel, Carrera de Ingeniería de Software, UTEQ.

\subsection*{Declaración de uso de asistentes de inteligencia artificial}

\begin{table}[H]
\centering
\small
\begin{tabular}{p{3cm}p{2.5cm}p{7.5cm}}
\toprule
\textbf{Herramienta} & \textbf{Versión/modelo} & \textbf{Propósito y sección donde se usó} \\
\midrule
Claude (Anthropic), vía Claude Code & Claude Sonnet 5 & Implementación de código (backend, frontend,
infraestructura), redacción y corrección de documentación técnica, ejecución y análisis de pruebas
empíricas (k6, Lighthouse, JaCoCo, OWASP ZAP, find-sec-bugs), búsqueda y verificación de referencias
bibliográficas, y redacción de este informe final completo (todas las secciones). Detalle completo en
\texttt{docs/etica/ai-disclosure.md}. \\
\bottomrule
\end{tabular}
\caption{Generative AI usage disclosure.}
\end{table}

Se declara explícitamente que ninguna cifra empírica de este informe fue generada por el asistente de IA
sin ejecutar realmente la herramienta correspondiente contra el sistema real --- ver la política completa
en \texttt{docs/etica/ai-disclosure.md}, que documenta también los episodios detectados de datos
fabricados en fases anteriores del proyecto (no atribuibles a este asistente, dado que ocurrieron antes
de su intervención documentada en el historial de commits) y cómo se corrigieron.

\subsection*{Agradecimientos}
Los autores agradecen al Ing. Guerrero Ulloa Gleiston Cíceron por la retroalimentación docente que
motivó las correcciones documentadas en \texttt{docs/observaciones/OBSERVACIONES.md}, y al Equipo E
(Tejada Bajaña Luis Alejandro, Alava Alvarado Jean Pierre) por la evaluación cruzada que identificó las
brechas E1--E10 cerradas en la Fase 5 de este proyecto.
\newpage
